\documentclass[11pt,a4paper]{article}

%% PACHETE MATEMATICE
\usepackage{amsmath,amssymb,amsfonts}
\usepackage{amsthm}
\usepackage{mathpazo}
\usepackage{eulervm}
\usepackage{enumerate}
\usepackage{color}
\usepackage{graphicx}
\usepackage[all]{xy}
\usepackage{xcolor}
\usepackage{framed}
\usepackage[unicode]{hyperref}
\setcounter{MaxMatrixCols}{20}
\usepackage{multicol}
%\usepackage[romanian]{babel}


%% ASPECT
\linespread{1.05}
\usepackage[T1]{fontenc}
\usepackage[utf8x]{inputenc}
\usepackage[margin=2cm]{geometry}
% \usepackage[color]{showkeys}
\usepackage{anyfontsize}
\usepackage{titlesec}
\titleformat*{\section}{\large\bfseries}

\usepackage{fancyhdr}
\pagestyle{fancy}
\rhead{\small \color{gray} Notițe de Adrian Manea}
\lhead{\small \color{gray} BCD-Aero, 2017---2018, L. Costache \& M. Olteanu}


\usepackage[autostyle = false, style = german]{csquotes}
\usepackage[romanian]{babel}
\newcommand\qq{\enquote}

%% COMENZILE MELE
\renewcommand*{\proofname}{Dem.:}
\newcommand\bb{\mathbb}
\newcommand\dr{\mathrm}
\newcommand\kal{\mathcal}
\newcommand\fk{\mathfrak}
\newcommand\op{\oplus}
\newcommand\ot{\otimes}
\newcommand\ds{\displaystyle}
\newcommand\id{\indent}
\newcommand\nid{\noindent}
\newcommand\seq{\subseteq}
\newcommand\sneq{\subsetneq}
\newcommand\speq{\supseteq}
\newcommand\spneq{\supsetneq}
\newcommand\rar{\rightarrow}
\newcommand\Rar{\Rightarrow}
\newcommand\xrar{\xrightarrow}
\newcommand\lar{\leftarrow}
\newcommand\Lar{\Leftarrow}
\newcommand\xlar{\xleftarrow}
\newcommand\lrar{\leftrightarrow}
\newcommand\Lrar{\Leftrightarrow}
\newcommand\ideal{\trianglelefteq}
\newcommand\ai{\text{ a.\^{i}. }}
\newcommand\sk{(\textit{Schi\c{t}\u{a}}) }
\newcommand\rhup{\rightharpoonup}
\newcommand\rhdn{\rightharpoondown}
\newcommand\lhup{\leftharpoonup}
\newcommand\lhdn{\leftharpoondown}
\newcommand\vid{\emptyset}
\newcommand\wh{\widehat}
\newcommand\ol{\overline}
\newcommand\NN{\mathbb{N}}
\newcommand\RR{\mathbb{R}}
\newcommand\ZZ{\mathbb{Z}}
\newcommand\QQ{\mathbb{Q}}
\newcommand\CC{\mathbb{C}}

%% STILURILE MELE
\newtheoremstyle{dotless}{}{}{\itshape}{}{\bfseries}{:}{ }{}

\theoremstyle{dotless}
\newtheorem{theorem}{Teorem\u{a}}[section]
\newtheorem{proposition}{Propozi\c{t}ie}[section]
\newtheorem{lemma}{Lem\u{a}}[section]
\newtheorem{corollary}{Corolar}[section]
\newtheorem{exercise}{Exerci\c{t}iu}

\newtheoremstyle{dotlessdr}{}{}{}{}{\bfseries}{:}{ }{}
\theoremstyle{dotlessdr}
\newtheorem{example}{Exemplu}[section]
\newtheorem{definition}{Defini\c{t}ie}[section]
\newtheorem{remark}{Observa\c{t}ie}[section]




\begin{document}

\begin{center}
  {
    \large\textbf{
      Seminar 9 \\
      Extremele funcțiilor de mai multe variabile \\
      Metoda celor mai mici pătrate.
    }}
\end{center}

\vspace{1cm}

\section{Extreme în mai multe dimensiuni}
\label{sec:extreme}

Pornim cu următoarea:

\begin{definition}\label{def:p-critic}
  Fie $f : A \to \RR$, cu $A \seq \RR^n$.

  Un punct $a \in A$ se numește \emph{punct critic} pentru $f$ dacă $f$ este
  diferențiabilă în $a$ și $df(a) = 0$, adică $\frac{\partial f}{\partial x_k} = 0$,
  pentru orice $k = 1, \dots, n$.
\end{definition}

Rezultă de aici că punctele de extrem local ale lui $f$ sînt printre soluțiile
sistemului:
\[
  \Big\{ \frac{\partial f}{\partial x_k}(x_1, \dots, x_n) = 0 \Big\}_k.
\]

În studiul naturii punctelor critice pentru o funcție $f$, pașii de urmat sînt:
\begin{enumerate}[(1)]
\item Se determină punctele critice, din anularea derivatelor parțiale;
\item Fie $a$ un punct critic. Se calculează matricea hessiană corespunzătoare,
  adică $H_f(a) = \Big( \frac{\partial^2 f}{\partial x_i \partial x_j} (a) \Big)_{i,j}$;
\item Se calculează valorile proprii ale lui $H$:
  \begin{enumerate}[(a)]
  \item Dacă toate valorile proprii sînt pozitive, $a$ este un minim local;
  \item Dacă toate sînt negative, $a$ este un maxim local;
  \item Dacă valorile proprii nu au semn uniform, $a$ nu este de extrem;
  \item Dacă 0 este valoare proprie, nu se poate decide. Astfel, se dezvoltă
    în serie Taylor a lui $f$ în jurul lui $a$, de unde se calculează semnul
    diferenței $f(x) - f(a)$.
  \end{enumerate}
\end{enumerate}

Pentru cazul bidimensional ($n = 2$), avem o metodă simplificată:
\begin{enumerate}[(1)]
\item Se determină punctele de extrem din sistemul:
  \[
    \frac{\partial f}{\partial x} = \frac{\partial f}{\partial y} = 0;
  \]
\item Fie $a = (a_1, a_2)$ un asemenea punct critic. Calculăm numerele:
  \begin{align*}
    r_0 &= \frac{\partial^2 f}{\partial x^2} (a_1, a_2) \\
    s_0 &= \frac{\partial^2 f}{\partial x\partial y}(a_1, a_2) \\
    t_0 &= \frac{\partial^2 f}{\partial y^2} (a_1, a_2).
  \end{align*}
  Observăm că avem:
  \[
    d^2 f(a_1, a_2) = r_0 dx^2 + 2s_0 dxdy + t_0 dy^2.
  \]
  \begin{enumerate}[(a)]
  \item Dacă $r_0 > 0$ și $r_0 t_0 - s_0^2 > 0$, atunci $a$ e punct de minim local;
  \item Dacă $r_0 < 0$ și $r_0 t_0 - s_0^2 > 0$, atunci $a$ e punct de maxim local;
  \item Dacă $r_0 t_0 < 0$, atunci $a$ nu este punct de extrem local;
  \item Dacă $r_0t_0 - s_0^2$, atunci studiem semnul $f(x_1, x_2) - f(a_1, a_2)$
    prin dezvoltare în serie Taylor.
  \end{enumerate}
\end{enumerate}

\textbf{Dezvoltarea în serie Taylor pentru o funcție de două variabile reale}
în jurul punctului $a = (x_0, y_0)$ este dată de formula:
\begin{align*}
  f(x, y) &= f(a) + \frac{1}{1!} \Big[ (x - x_0) \frac{\partial f}{\partial x}(a) + %
            (y - y_0) \frac{\partial f}{\partial y}(a) \Big]\\
          &+ \frac{1}{2!} \Big[ (x - x_0)^2 \frac{\partial^2 f}{\partial x^2}(a) + %
            2(x - x_0)(y - y_0) \frac{\partial^2 f}{\partial x \partial y}(a) + %
            (y - y_0)^2 \frac{\partial^2 f}{\partial y^2}(a) \Big] + \\
          &+ \frac{1}{3!} \Big[ (x - x_0)^3 \frac{\partial^3 f}{\partial x^3} (a) + %
            3(x - x_0)^2(y - y_0) \frac{\partial^3 f}{\partial x^2 \partial y}(a) + %
            3(x - x_0)(y - y_0)^3 \frac{\partial^3 f}{\partial x \partial y^2}(a) + %
            (y - y_0)^3 \frac{\partial^3 f}{\partial y^3}(a) \\
          &+ \dots
\end{align*}

\begin{remark}\label{rk:frontiera}
  În cazul în care se impun constrîngeri pentru domeniul de definiție a funcției,
  se studiază separat problema punctelor de extrem în interiorul domeniului,
  precum și pe frontieră.
\end{remark}


%%%%%%%%%%%%%%%%%%%%%%%%%%%%%%%%%%%%%%%%%%%%%%%%%%%%%%%%%%%%%%%%%%%%%%


\section{Regresia liniară}
\label{sec:reg-lin}

Regresia liniară se adresează unor probleme de tip statistic, atunci cînd avem
la dispoziție o serie de date experimentale, pe care trebuie să le corelăm cu
un model teoretic \emph{liniar}.

Așadar, date o serie de perechi de puncte de forma $(x_i, y_i)$, care reprezintă
datele colectate într-un experiment care verifică o teorie bazată pe ecuații
liniare, ne punem problema găsirii celei mai potrivite drepte care să prezinte
o abatere minimă față de punctele colectate (eng. ``best fit'').

Fie, așadar, dreapta $y = f(x) = ax + b$ modelul teoretic pe care îl studiem.
Dacă $(x_i, y_i)$ reprezintă colecția de puncte experimentale, atunci vom minimiza
abaterea lor de la funcția $f$ prin intermediul \emph{metodei celor mai mici pătrate}.
Vom impune, așadar, ca suma pătratelor distanțelor de la punctele colectate la
dreapta teoretică să fie minimă.

În expresia $f(x) = ax + b$, necunoscutele sînt $a$ și $b$, deoarece ele determină
dreapta experimentală pe care o căutăm. Așadar, putem privi $f$ ca pe o funcție
de variabilele $a$ și $b$. Scriem, deci, expresia corespunzătoare celor mai
mici pătrate:
\[
  S(a, b) = \sum_{i = 1}^n [ f(x_i) - y_i ]^2.
\]

Această funcție se dorește minimizată, așadar problema revine la găsirea punctelor
de minim pentru ea.

Se studiază punctele critice cu metoda din secțiunea de mai sus, apoi se
găsesc punctele $(a, b)$ care să fie de minim.

În fine, \emph{dreapta de regresie} este $y = ax + b$.

Se poate face și \emph{calculul erorilor}, folosind variabilele:
\begin{align*}
  S_r &= \sum_i (y_i - b - ax_i)^2 \\
  S_t &= \sum_i (y_i - \overline{y})^2,
\end{align*}
unde $\overline{y}$ este media (aritmetică) a valorilor $y_i$ experimentale.

În fine, se pot calcula \emph{eroarea standard} a estimării și, respectiv,
abaterea medie pătratică:
\[
  \varepsilon = \sqrt{\frac{S_r}{n - 2}}, \quad \sigma = \sqrt{\frac{S_t}{n - 1}}.
\]

%%%%%%%%%%%%%%%%%%%%%%%%%%%%%%%%%%%%%%%%%%%%%%%%%%%%%%%%%%%%%%%%%%%%%%

\section{Exerciții}
\label{sec:exc}

\indent\indent 1. Fie funcția:
\[
  f : D \seq \RR^2 \to \RR, f(x, y) = x^2 + xy + y^2 - 4\ln x - 10 \ln y + 3.
\]
Determinați punctele de extrem și calculați valorile funcției în aceste puncte.

\vspace{1cm}

2. Fie $f : D \seq \RR^2 \to \RR, f(x, y) = x^3 + y^3 - 6xy$.
\begin{enumerate}[(a)]
\item Pentru $D = \RR^2$, determinați punctele de extrem și calculați valorile
  funcției în aceste puncte;
\item Pentru $D = \{ (x,y) \in \RR^2 \mid x \geq 0, y \geq 0, x + y \geq 5 \}$,
  determinați valoarea minimă și maximă a funcției.
\end{enumerate}

Indicație (b): Considerați funcțiile $g_1(x) = f(x,0), g_2 = f(0, y)$ și apoi
funcția $g_3 = f(x, 5-x)$, cărora le găsiți punctele de extrem.

\vspace{1cm}

3. Fie $f : D \to \RR^2, f(x, y) = 3xy^2 - x^3 - 15x - 36y + 9$.
\begin{enumerate}[(a)]
\item Pentru $D = \RR^2$, determinați punctele de extrem și calculați valorile
  funcției în aceste puncte;
\item Pentru $D = [-4, 4] \times [-3, 3]$ determinați valoarea minimă și valoarea
  maximă a funcției.
\end{enumerate}

Indicație (b): Considerați funcțiile $f(x, -3), f(4, y), f(x, 3), f(-4, y)$.

\vspace{1cm}

4. Fie $f : D \seq \RR^2 \to \RR, f(x, y) = 4xy - x^4 - y^4$.
\begin{enumerate}[(a)]
\item Pentru $D = \RR^2$, determinați punctele de extrem și calculați valorile
  funcției în aceste puncte;
\item Pentru $D = [-1, 2] \times [0, 2]$, determinați valoarea minimă și
  maximă a funcției.
\end{enumerate}

\vspace{1cm}

5. Fie $f : D \seq \RR^2 \to \RR, f(x, y) = x^3 + 3x^2 y - 15x - 12 y$.
\begin{enumerate}[(a)]
\item Pentru $D = \RR^2$, determinați punctele de extrem și calculați
  valorile funcției în aceste puncte;
\item Pentru $D = \{ (x, y) \in \RR^2 \mid x \geq 0, y \geq 0, 3y + x \leq 3 \}$
  determinați valoarea minimă și valoarea maximă a funcției.
\end{enumerate}

\vspace{1cm}

6. Determinați valorile extreme pentru funcțiile $f$, definite pe domeniile $D$,
unde:
\begin{enumerate}[(a)]
\item $f(x, y) = x^3 + y^3 - 6xy, D = \RR^2$;
\item $f(x, y) = xy(1 - x - y), D = [0, 1] \times [0, 1]$;
\item $f(x, y) = x^4 + y^4 - 2x^2 + 4xy - 2y^2, D = (-\infty, 0) \times (0, \infty)$;
\item $f(x, y) = x^3 + 8y^3 - 2xy, D = \{(x, y) \in \RR^2 \mid x, y \geq 0, y + 2x \leq 2\}$;
\item $f(x, y) = x^4 + y^3 - 4x^3 - 3y^2 + 3y, D = \{ (x, y) \in \RR^2 \mid x^2 + y^2 < 4 \}$.
\end{enumerate}

\vspace{1cm}

7. Să se determine dreapta de regresie care mediază printre punctele
$M_1(1, 2), M_2(2, 0), M_3(3, 1)$.

\vspace{1cm}

8*. Dintre toate paralelipipedele dreptunghice cu volum constant 1, determinați
pe cel cu aria totală minimă.

\end{document}
